\documentclass[aps,preprint]{revtex4}%
\usepackage{amssymb}
\usepackage{amsfonts}%
\usepackage{amsmath}%
\setcounter{MaxMatrixCols}{30}%
\usepackage{graphicx}
%TCIDATA{OutputFilter=latex2.dll}
%TCIDATA{Version=5.50.0.2960}
%TCIDATA{CSTFile=revtex4.cst}
%TCIDATA{Created=Sunday, October 21, 2018 07:48:13}
%TCIDATA{LastRevised=Monday, October 22, 2018 09:47:45}
%TCIDATA{<META NAME="GraphicsSave" CONTENT="32">}
%TCIDATA{<META NAME="SaveForMode" CONTENT="1">}
%TCIDATA{BibliographyScheme=Manual}
%TCIDATA{<META NAME="DocumentShell" CONTENT="Articles\SW\REVTeX 4">}
%BeginMSIPreambleData
\providecommand{\U}[1]{\protect\rule{.1in}{.1in}}
%EndMSIPreambleData
\newtheorem{theorem}{Theorem}
\newtheorem{acknowledgement}[theorem]{Acknowledgement}
\newtheorem{algorithm}[theorem]{Algorithm}
\newtheorem{axiom}[theorem]{Axiom}
\newtheorem{claim}[theorem]{Claim}
\newtheorem{conclusion}[theorem]{Conclusion}
\newtheorem{condition}[theorem]{Condition}
\newtheorem{conjecture}[theorem]{Conjecture}
\newtheorem{corollary}[theorem]{Corollary}
\newtheorem{criterion}[theorem]{Criterion}
\newtheorem{definition}[theorem]{Definition}
\newtheorem{example}[theorem]{Example}
\newtheorem{exercise}[theorem]{Exercise}
\newtheorem{lemma}[theorem]{Lemma}
\newtheorem{notation}[theorem]{Notation}
\newtheorem{problem}[theorem]{Problem}
\newtheorem{proposition}[theorem]{Proposition}
\newtheorem{remark}[theorem]{Remark}
\newtheorem{solution}[theorem]{Solution}
\newtheorem{summary}[theorem]{Summary}
\newenvironment{proof}[1][Proof]{\noindent\textbf{#1.} }{\ \rule{0.5em}{0.5em}}
\begin{document}
\preprint{HEP/123-qed}
\title[Short title for running header]{REV\TeX4}
\author{First Author}
\affiliation{My Institution}
\author{Second Author}
\affiliation{My Institution}
\author{Third Author}
\affiliation{Other Institution}
\keywords{one two three}
\pacs{PACS number}

\begin{abstract}
Shell document for REV\TeX{} 4.

\end{abstract}
\volumeyear{year}
\volumenumber{number}
\issuenumber{number}
\eid{identifier}
\date[Date text]{date}
\received[Received text]{date}

\revised[Revised text]{date}

\accepted[Accepted text]{date}

\published[Published text]{date}

\startpage{101}
\endpage{102}
\maketitle
\tableofcontents


\section{Operator Systems}

Kadison proved in [5] that every ordered real vector space with an Archimedean
order unit can be represented as a vector subspace of the space of continuous
real-valued functions on a compact Hausdorff space via an order-preserving map
that carries the order unit to the constant function $1$. Such subspaces are
generally referred to as \emph{function systems}, and Kadison's result
therefore provides a representation theorem for real function systems.
Kadison's representation theorem motivated Choi and Effros [2] to obtain an
analogous representation theorem for self-adjoint subspaces of unital
$c^{\ast}$- algebras that contain the unit. Such subspaces, together with a
natural matrix order, are called \emph{operator systems}. Operator systems are
complex vector spaces with a $\ast$-operation, and we refer to these as $\ast
$-vector spaces.

\begin{enumerate}
\item[a] In any vector space $V$ , the self-adjoint elements $V_{h}$ form a
real subspace and one may define an ordering on $V$ by specifying a
\emph{cone} of positive elements $V_{+}\subset V_{h}$. When $\left\{
V,V_{+}\right\}  $ is an ordered vector space, then $\left\{  V_{h}%
,V_{+}\right\}  $ is an ordered real vector space. In addition, an element
$e\in V$ is an \emph{order unit} (respectively, an \emph{Archimedean order}
unit) for $\left\{  V,V_{+}\right\}  $ if it is an order unit (respectively,
an Archimedean order unit) for $\left\{  V_{h},V_{+}\right\}  $. The axioms
for an operator system require, among other things, that it is an ordered
$\ast$-vector space with Archimedean order unit $e$ and that for each
$n\in\mathbb{N}$ the $\ast$-vector space $M_{n}(V)$ is ordered with
Archimedean order unit
\begin{equation}
e_{n}=\left(
\begin{array}
[c]{cccc}%
e & 0 & \cdots & 0\\
0 & \ddots &  & 0\\
\vdots &  & \ddots & \vdots\\
0 & 0 &  & 0
\end{array}
\right)
\end{equation}

\end{enumerate}

\subsection{Cones}

\begin{definition}
If V is a real vector space, a cone in V is a nonempty subset $C\subseteq V$
with the following two properties:
\end{definition}

\begin{enumerate}
\item $\alpha v$ $\in\mathcal{C}$ whenever $\alpha\in\lbrack0,\infty)$ and $v$
$\in\mathcal{C}$

\item $v+w\in\mathcal{C}$ whenever $v,w\in\mathcal{C}$.

\emph{An ordered vector space} $\left\{  V,V_{+}\right\}  $ is a pair
consisting of a real vector space $V$ and a cone $V_{+}\subseteq V$ satisfying

\item  $V_{+}\cap\left(  -V_{+}\right)  $ $=\{0\}$.
\end{enumerate}

In any ordered vector space we may define a \emph{partial ordering} (i.e.,
\emph{a reflexive, antisymmetric, and transitive relation}) on $V$ by defining
$v\geq w$ if and only if $v-w\in V_{+}$. Note that this partial ordering is
$\emph{translation}$ $\emph{invariant}$ (i.e., $v\geq w$ implies $v+x\geq
w+x$) and invariant under multiplication by non-negative reals (i.e., $v\geq
w$ and $\alpha\in[0,\infty)$ implies $\alpha v\geq\alpha w$). Also note that
$v\in V_{+}$ if and only if $v\geq0$; for this reason $V_{+}$ is sometimes
called the cone of positive elements of $V$ .

\begin{remark}
Although we have used the set $V_{+}$ to define a partial ordering, we could
just as easily have gone the other direction: If is a partial ordering on $V$
that is translation invariant and invariant under multiplication by
non-negative reals, then the set $V_{+}$ $:=v\in V_{+}:v\geq0$ is a set
satisfying $(1),(2)$ and $(3)$ above, and consequently $\left\{
V,V_{+}\right\}  $ is an ordered vector space. 
\end{remark}

\begin{remark}
There is some disagreement in the literature about the use of the term
\textquotedblleft\emph{cone}\textquotedblright. While there are authors that
conform with the definitions as we have given them there are other authors who
use the term \textquotedblleft\emph{cone}\textquotedblright\ to refer to a
subset of V satisfying $(1),(2),$ and $(3)$ stated above, and use the term
\textquotedblleft\emph{wedge}\textquotedblright\ for a subset satisfying only
$(1)$ and $(2)$.
\end{remark}

\begin{definition}
If $\left\{  V,V_{+}\right\}  $ is an ordered real vector space, an element
$e\in V$ is called an \emph{order unit} for $V$ if for each $v\in V$ there
exists a real number $r>0$ such that $re\geq v$.
\end{definition}

\begin{definition}
If $\left\{  V,V_{+}\right\}  $ is an ordered real vector space, then we say that
\end{definition}

$V_{+}$ is \emph{full} (or that $V_{+}$ is a full cone) if $V=V_{+}-V_{+}$.

\begin{lemma}
If $\left\{  V,V_{+}\right\}  $ is an ordered real vector space with order
unit $e$, then
\end{lemma}

\begin{enumerate}
\item $e\in V_{+}$

\item If $v\in V$ and a real number $r>0$ is chosen so that $re\geq v$, then
$se\geq v$ for all $s\geq r$

\item If $v_{1},...,v_{n}\in V$, then there exists $r>0$ such that $re\geq
v_{i}$ for all $1\leq i\leq n$

\item $V_{+}$ is a full cone of $V$ 

\item If $v_{1},...,v_{n}\in V_{+}$ and $v_{1}+...+v_{n}=0$, then
$v_{1}=...=v_{n}=0$ 

\item Suppose $v_{1},...,v_{n}\in V_{+}$ and $\alpha_{1},...,\alpha_{n}$
$\in\lbrack0,\infty)$. If $\alpha_{1}v_{1}+...+\alpha_{1}v_{n}=0$, then for
each $1\leq i\leq n$ either $\alpha_{i}=0$ or $v_{i}=0$. 
\end{enumerate}

\begin{definition}
If $\left\{  V,V_{+}\right\}  $ is an ordered real vector space with an order
unit $e$, then we say that $e$ is an Archimedean order unit if whenever $v\in
V$ with $re+v\geq0$ for all real $r>0$, then $v\in V_{+}$.
\end{definition}

\begin{lemma}
Let $\left\{  V,V_{+}\right\}  $ be an ordered real vector space with an
Archimedean order unit $e$, and let $r_{0}\in\lbrack0,\infty)$. If $v\in V$
and $re+v\geq0$ for all $r>r_{0}$, then $r_{0}e+v\geq0$.
\end{lemma}

\subsection{Positive $\mathbb{R}$-Linear Functionals and States on Real
Spaces}

\begin{definition}
If $\left\{  V,V_{+}\right\}  $ is an ordered real vector space, an
$\mathbb{R}$-linear functional $f:V\rightarrow R$ is called positive if
$f(V_{+})$ $\subseteq\lbrack0,\infty)$.
\end{definition}

\begin{definition}
If $\left\{  V,V_{+}\right\}  $ is an ordered real vector space and
$S\subseteq V$ , we say that $S$ \emph{majorizes} $V_{+}$ if for each $v\in
V_{+}$ there exists $w\in S$ such that $w\geq v$.
\end{definition}

\begin{remark}
Note that if $S$ majorizes $V_{+}$ and $S\subseteq T$ , then $T$ majorizes
$V_{+}$. In addition, if $e$ is an order unit for $\left\{  V,V_{+}\right\}
$, and if $E$ \emph{is a subspace of} $V$ containing $e$, then $E$ majorizes
$V_{+}$. This is due to the fact that if $v\in V$ , then there exists a real
$r>0$ such that $re\geq v$, and since $E$ is a subspace we have that $re\in E$.
\end{remark}

The next theorem may be thought of \emph{as an analogue of the Hahn-Banach
Theorem} since it gives conditions under which we may \emph{extend a positive
}$\mathbb{R}$\emph{-linear functional} on a subspace of $V$ to \emph{a
positive }$\mathbb{R}$\emph{-linear functional} on all of $V$ . 

\begin{theorem}
Suppose $\left\{  V,V_{+}\right\}  $ is an ordered real vector space and that
$V_{+}$ is a full cone for $V$ . If $E$ is a subspace of $V$ that majorizes
$V_{+},$ and if $f:E\rightarrow\mathbb{R}$ is a positive $\mathbb{R}$-linear
functional on $E$, then there exists a positive $\mathbb{R}$-linear functional
$g:V\rightarrow\mathbb{R}$ such that $\left.  g\right\vert _{E}=f$.
\end{theorem}

\subsection{Positive $\mathbb{C}$-Linear Functionals and States on Complex
Spaces}

In this section we will consider complex vector spaces with a $\ast
$-operation, and describe what it means for these spaces to be ordered. We
also develop analogues of the results in the last subsection for these ordered
complex vector spaces.

\begin{definition}
A $\ast$-vector space consists of a complex vector space $V$ together with a
map $\ast$ $:V\rightarrow V$ that is involutive (i.e.,$(v^{\ast})^{\ast}=v$
for all $v\in V$) and conjugate linear (i.e., $(\U{3bb} v+w)^{\ast}=\U{3bb}
v^{\ast}+w^{\ast}$ for all $\U{3bb} \in\mathbb{C}$ and $v,w$ $\in V$ ). If $V$
is a $\ast$-vector space, then we let $V_{h}:=\{x\in V:x^{\ast}=x\}$ and we
call the elements of $V_{h}$ the hermitian elements of $V$ .
\end{definition}

It is easy to see that $V_{h}$ is a real subspace of the complex vector space
$V$ . Also note that any $v\in V$ may be written uniquely as $v=x+iy$ with
$x,y\in$ $V_{h}$. In fact, $v=x+iy$ with $x,y\in V_{h}$ if and only if
$x=(v+v^{\ast})/2$ and $y=(v-v^{\ast})/2i$. We call $x$ and $y$ the real part
and imaginary part of $v$, respectively, and we write%
\begin{align}
\operatorname{Re}\left\{  v\right\}    & =(v+v^{\ast})/2\\
\operatorname{Im}\left\{  v\right\}    & =(v-v^{\ast})/2i
\end{align}




Note that we also have
\begin{equation}
V\cong V_{h}\oplus iV_{h}%
\end{equation}
as real vector spaces.

\begin{definition}
If $V$ is a $\ast$-vector space, we say that $\left\{  V,V_{+}\right\}  $ is
an ordered $\ast$-vector space if $V_{+}\subseteq V_{h}$ and the following two
conditions hold:
\end{definition}

\begin{enumerate}
\item $V_{+}$ is a cone of $V_{h}$ (i.e., $V_{+}+V_{+}\subseteq V_{+}$ and
$aV_{+}\subseteq V_{+}$ for all $a\in\lbrack0,\infty))$.

\item $V_{+}\cap-V_{+}=\{0\}$.

In this case for $v,w\in V_{h}$ we write $v\geq w$ to mean that $v-w\in V_{+}$.
\end{enumerate}

\begin{definition}
If $\left\{  V,V_{+}\right\}  $ is an ordered $\ast$-vector space, then an
element $e\in$ $V_{h}$ is an order unit if for all $x\in V_{h}$ there exists a
real number $r>0$ such that $re\geq x$. If $e$ is an order unit, we say that
$e$ is Archimedean if whenever $x\in V_{h}$ and $re+x\geq0$ for all
$r\in(0,\infty)$, then $x\geq0$.
\end{definition}

\begin{remark}
. Note that $\left\{  V,V_{+}\right\}  $ is an ordered -vector space if and
only if $\left\{  V_{h},V_{+}\right\}  $ is an ordered real vector space.
Similarly, $e$ is an order unit for the ordered vector space $\left\{
V,V_{+}\right\}  $ if and only if $e$ is an order unit for the ordered real
vector space $\left\{  V_{h},V_{+}\right\}  $; and $e$ is Archimedean for
$\left\{  V,V_{+}\right\}  $ if and only if $e$ is Archimedean for $\left\{
V_{h},V_{+}\right\}  .$  
\end{remark}

We now turn our attention to positive maps, positive $\mathbb{C}$-linear
functionals, and states on ordered $\ast$-vector spaces. Many of the results
here may be viewed as analogues of the previous section.

\begin{definition}
Let $\left\{  V,V_{+}\right\}  $ be an ordered -vector space with order unit
$e$. A $\mathbb{C}$-linear functional $f:V\rightarrow\mathbb{C}$ is positive
if $f(V_{+})\subset\lbrack0,)$ and a state if it is positive and $f(e)=1$.
\end{definition}

\begin{definition}
. Let $\left\{  V,V_{+}\right\}  $ be an ordered $\ast$-vector space with
order unit $e$, and let $\left\{  W,W_{+}\right\}  $ be an ordered $\ast
$-vector space with order unit $e\prime$. A linear map $\U{3c6} :V\rightarrow
W$ is positive if $v\in V_{+}\Rightarrow\U{3c6} (v)\in W_{+}$, and unital if
$\U{3c6} (e)=e\prime$. A linear map $\U{3c6} :V\rightarrow W$ is an order
isomorphism if $\U{3c6} $ is bijective and $v\in V_{+}$ if and only if
$\U{3c6} (v)\in W_{+}.$
\end{definition}

\begin{proposition}
Let $\left\{  V,V_{+}\right\}  $ be an ordered $\ast$-vector space with an
order unit, and let $\left\{  W,W_{+}\right\}  $ be an ordered $\ast$-vector
space. If $\U{3c6} :V\rightarrow W$ is a positive linear map, then $\U{3c6}
(v^{\ast})=\U{3c6} (v)^{\ast}$for all $v\in V$ .
\end{proposition}

\begin{corollary}
If $\left\{  V,V_{+}\right\}  $ is an ordered $\ast$-vector space with an
order unit and $f:V\rightarrow\mathbb{C}$ is a positive linear functional,
then $f(v^{\ast})=f(v)^{\ast}$ for all $v\in V$ .
\end{corollary}

\begin{definition}
Let $\left\{  V,V_{+}\right\}  $ be an ordered $\ast$-vector space. If
$f:V_{h}\rightarrow\mathbb{R}$, then we define $\widetilde{f}:V\rightarrow
\mathbb{C}$ by
\end{definition}

$\widetilde{f}(v):=f(\operatorname{Re}(v))+if(\operatorname{Im}(v))$.

\begin{proposition}
Let $\left\{  V,V_{+}\right\}  $ be an ordered $\ast$-vector space. If
$f:V_{h}\rightarrow\mathbb{R}$ is R-linear, then $\widetilde{f}:V\rightarrow
\mathbb{C}$ is $\mathbb{C}$-linear. In addition, $\widetilde{f}$ is positive
if and only if $f$ is positive, and $\widetilde{f}$ is a state if and only if
$f$ is a state.
\end{proposition}

The following proposition allows us to characterize the positive linear
functionals on ordered $\ast$-vector spaces with an Archimedean unit.

\begin{proposition}
Let $\left\{  V,V_{+}\right\}  $ be an ordered $\ast$-vector space with order
unit $e$. If $f:V\rightarrow\mathbb{C}$ is a $\mathbb{C}$-linear functional,
then $f$ is positive if and only if $f=g$ for some positive $\mathbb{R}%
$-linear functional $g:V_{h}\rightarrow\mathbb{R}$. 
\end{proposition}

\begin{proposition}
Let $\left\{  V,V_{+}\right\}  $ be an ordered -vector space with Archimedean
order unit $e$. If $v\in V$ and $f(v)=0$ for every state $f:V\rightarrow
\mathbb{C}$, then $v=0$.
\end{proposition}

\begin{proposition}
Let $\left\{  V,V_{+}\right\}  $ be an ordered $\ast$-vector space with
Archimedean order unit $e$. If $v\in V$ and $f(v)\geq0$ for every state
$f:V\rightarrow\mathbb{C}$, then $v\in V_{+}$.
\end{proposition}

\begin{remark}
Note that when $V$ is a $\ast$-vector space, we have $V=V_{h}\oplus iV_{h}$.
In fact, $V$ is the complexification of the real vector space $V_{h}$, and
alternatively we may describe $V$ as the complex vector space $V_{h}%
\otimes_{\mathbb{R}}\mathbb{C}$. Likewise, if $f:V_{h}\rightarrow\mathbb{R}$
is $\mathbb{R}$-linear, the $\mathbb{C}$-linear map $\widetilde{f}%
:V\rightarrow\mathbb{C}$ is the complexification of $f$ . The complexification
is an (additive) \emph{functor} from the \emph{category of real vector spaces
to the category of complex vector spaces}.
\end{remark}

\begin{theorem}
Suppose $\left\{  V,V_{+}\right\}  $ is an ordered $\ast$-vector space and
that $V_{+}$ is a full cone for $V$ . If $E$ is a subspace of $V$ that
majorizes $V_{+},$ and if $f:E\rightarrow\mathbb{C}$ is a positive
$\mathbb{C}$-linear functional on $E$, then there exists a positive
$\mathbb{C}$-linear functional $g:V\rightarrow\mathbb{C}$ such that $\left.
g\right\vert _{E}=f$.
\end{theorem}

Examples of $\ast$-vector spaces with Archimedean order unit $e$ is are unital
$c^{\ast}$-algebras e.g. let $\mathcal{A}\subseteq\mathcal{B}(\mathcal{H})$ be
a unital $c^{\ast}$-algebras with unit $e=I_{\mathcal{H}}$, then $\mathcal{A}$
is an Archimedean ordered $\ast$-vector space with $\mathcal{A}_{+}$ equal to
the usual cone of positive elements. 


\end{document}